\documentclass{article}
\usepackage{amsmath,amssymb,amsthm}
\usepackage{algorithm,algorithmic}
\usepackage{hyperref}
\usepackage{enumitem}
\usepackage{bm}
\newtheorem{theorem}{Theorem}
\newtheorem{lemma}{Lemma}
\newtheorem{assumption}{Assumption}
\newcommand{\E}{\mathbb{E}}
\newcommand{\R}{\mathbb{R}}

\begin{document}

\section*{AdaGrad with Scalar Adaptive Stepsize}

The algorithm is the scalar version of AdaGrad: it accumulates the squared gradients and scales the learning rate by the inverse square‑root of this accumulation.  The analysis below is carried out for a single scalar parameter $w\in\R$; the vector case follows by applying the same argument coordinate‑wise.

---------------------------------------------------------------------

\section*{Mathematical Formulation}

Let $f:\R\to\R$ be differentiable and denote the (stochastic) gradient at iteration $k$ by
\[
g_k \;=\; \nabla f(w_k;\,\xi_k),
\]
where $\xi_k$ is a random data sample (or minibatch).  
Define the accumulated squared gradient
\[
G_k \;:=\; \sum_{i=1}^{k} g_i^{\,2}\qquad (G_0:=0).
\]
Given a base stepsize $\alpha>0$ and a small constant $\varepsilon>0$ for numerical stability, the AdaGrad update reads
\begin{equation}
\boxed{
w_{k+1}\;=\;w_k \;-\; \frac{\alpha}{\sqrt{G_k+\varepsilon}}\; g_k.}
\tag{1}
\end{equation}
The algorithm keeps a running sum of the raw gradients
\[
A_k:=\sum_{i=1}^{k} g_i,
\]
which is not directly used in the update but can be useful for diagnostics.

---------------------------------------------------------------------

\section*{Pseudocode}

\begin{algorithm}[H]
\caption{Scalar AdaGrad}
\begin{algorithmic}[1]
\REQUIRE Initial point $w_0\in\R$, base stepsize $\alpha>0$, stability constant $\varepsilon>0$.
\STATE $G_0\gets 0$.
\FOR{$k=0,1,\dots,T-1$}
    \STATE Sample $\xi_k$ and compute stochastic gradient $g_k=\nabla f(w_k;\xi_k)$.
    \STATE $G_{k+1}\gets G_k+g_k^{\,2}$.
    \STATE $w_{k+1}\gets w_k-\dfrac{\alpha}{\sqrt{G_{k+1}+\varepsilon}}\;g_k$.
\ENDFOR
\RETURN $w_T$ (or the averaged iterate $\bar w_T:=\frac1T\sum_{k=1}^{T}w_k$).
\end{algorithmic}
\end{algorithm}

---------------------------------------------------------------------

\section*{Dry Run Example}

Consider the simple quadratic
\[
f(w)=(w-3)^2,
\]
so that $\nabla f(w)=2(w-3)$.  Use a deterministic gradient (i.e. $\xi_k$ is irrelevant), base stepsize $\alpha=0.1$, $\varepsilon=10^{-8}$ and start from $w_0=0$.

\begin{center}
\begin{tabular}{c|c|c|c|c}
Iteration $k$ & $w_k$ & $g_k=2(w_k-3)$ & $G_k$ & $w_{k+1}=w_k-\dfrac{0.1}{\sqrt{G_k+\varepsilon}}g_k$\\\hline
0 & $0$ & $-6$ & $0$ & $0-\dfrac{0.1}{\sqrt{0+10^{-8}}}(-6)=0+0.1\cdot6\cdot10^{4}=6000$\\
1 & $6000$ & $2(6000-3)=11994$ & $(-6)^2=36$ & $6000-\dfrac{0.1}{\sqrt{36+10^{-8}}}(11994)
=6000-\dfrac{0.1}{6}\,11994\approx6000-199.9\approx5800.1$\\
2 & $5800.1$ & $2(5800.1-3)=11594.2$ & $36+11994^{2}=36+1.4395\times10^{8}\approx1.4395\times10^{8}$ &
\[
w_{3}=5800.1-\frac{0.1}{\sqrt{1.4395\times10^{8}}}\,11594.2
\approx5800.1-\frac{0.1}{11996}\,11594.2\approx5800.1-0.967\approx5799.13.
\]
\end{tabular}
\end{center}
The objective values $f(w_k)$ decrease dramatically after the first gigantic step (caused by the tiny denominator) and then settle around the optimum $w^\star=3$.  In practice one chooses a much smaller $\alpha$ or a larger $\varepsilon$ to avoid such extreme behavior; the example merely illustrates the mechanics of the update.

---------------------------------------------------------------------

\section*{Assumptions}

\begin{assumption}[Smoothness]\label{ass:smooth}
The stochastic loss $f(\cdot;\xi)$ is $L$‑smooth in expectation, i.e.
\[
\forall w,w'\in\R:\qquad 
\E_{\xi}\!\big[\|\nabla f(w;\xi)-\nabla f(w';\xi)\|^2\big]\le L^2\|w-w'\|^2 .
\]
Equivalently, for the full objective $F(w):=\E_{\xi}[f(w;\xi)]$ we have
\[
F(w')\le F(w)+\langle\nabla F(w),w'-w\rangle+\frac{L}{2}\|w'-w\|^2 .
\tag{2}
\]
\end{assumption}

\begin{assumption}[Unbiased Gradient and Bounded Variance]\label{ass:unbiased}
For every $k$,
\[
\E[g_k\,|\,\mathcal{F}_k]=\nabla F(w_k),\qquad 
\E\big[\|g_k-\nabla F(w_k)\|^2\;\big|\;\mathcal{F}_k\big]\le\sigma^2,
\]
where $\mathcal{F}_k:=\sigma(w_0,\xi_0,\dots,\xi_{k-1})$ is the filtration generated by the algorithm up to iteration $k$.
\end{assumption}

\begin{assumption}[Bounded Gradients]\label{ass:bound}
There exists $G_{\max}>0$ such that $\|g_k\|\le G_{\max}$ almost surely for all $k$.
\end{assumption}

All constants $\alpha,\,\varepsilon,\,L,\,\sigma,\,G_{\max}$ are fixed and independent of $T$.

---------------------------------------------------------------------

\section*{Main Convergence Theorem}

\begin{theorem}[Average Gradient Norm of Scalar AdaGrad]\label{thm:main}
Under Assumptions~\ref{ass:smooth}–\ref{ass:bound}, for any $T\ge 1$ the iterates generated by \eqref{1} satisfy
\[
\frac{1}{T}\sum_{k=0}^{T-1}\E\!\big[\|\nabla F(w_k)\|^2\big]
\;\le\;
\underbrace{\frac{2(F(w_0)-F^\star)}{\alpha\,\sqrt{\varepsilon}}}_{\displaystyle C_1}
\;+\;
\underbrace{\frac{L\alpha G_{\max}}{\sqrt{\varepsilon}}}_{\displaystyle C_2}
\;+\;
\underbrace{\frac{L\alpha\sigma}{\sqrt{T\,\varepsilon}}}_{\displaystyle C_3},
\tag{3}
\]
where $F^\star:=\inf_{w}F(w)$.  Consequently,
\[
\min_{0\le k<T}\E\!\big[\|\nabla F(w_k)\|^2\big]=O\!\Big(\frac{1}{\sqrt{T}}\Big).
\]
\end{theorem}

The bound consists of a deterministic term $C_1$ (initial sub‑optimality), a term $C_2$ that captures the effect of the adaptive stepsize, and a stochastic term $C_3$ that decays as $1/\sqrt{T}$.

---------------------------------------------------------------------

\section*{Theoretical Properties}

\begin{enumerate}[label=\arabic*.]
\item \textbf{Stability.}  Because the denominator $\sqrt{G_k+\varepsilon}$ is non‑decreasing, the effective stepsize $\eta_k:=\alpha/\sqrt{G_k+\varepsilon}$ is non‑increasing.  Hence the algorithm cannot ``explode’’ after a large gradient; it automatically dampens subsequent updates.

\item \textbf{Hyper‑parameter Sensitivity.}
    \begin{itemize}
        \item $\alpha$ appears linearly in the constants $C_1$ and $C_2$; a larger $\alpha$ speeds up early progress (reduces $C_1$) but increases $C_2$.
        \item $\varepsilon$ controls the initial stepsize: $\eta_0=\alpha/\sqrt{\varepsilon}$.  Choosing $\varepsilon$ too small leads to excessively large first steps (as seen in the dry run), while a too‑large $\varepsilon$ makes AdaGrad behave like vanilla SGD.
    \end{itemize}
\item \textbf{Comparison to Baselines.}
    \begin{itemize}
        \item For standard (non‑adaptive) SGD with constant stepsize $\eta$, the best known guarantee for non‑convex smooth $F$ is $\mathcal{O}(1/\sqrt{T})$ on the average gradient norm, but it requires a careful tuning of $\eta$ that depends on $L$ and $G_{\max}$.
        \item AdaGrad removes the need to know $L$ or $G_{\max}$ a priori: the adaptivity automatically scales the stepsize, and the same $\mathcal{O}(1/\sqrt{T})$ rate is achieved with the universal choice $\alpha=1$ (up to constants).
    \end{itemize}
\end{enumerate}

---------------------------------------------------------------------

\section*{Preliminaries (Definitions and Lemmas)}

\begin{lemma}[Smoothness Inequality]\label{lem:smooth}
If $F$ is $L$‑smooth then for any $w,w'$
\[
F(w')\le F(w)+\langle\nabla F(w),w'-w\rangle+\frac{L}{2}\|w'-w\|^2 .
\]
\end{lemma}
\begin{proof}
Standard consequence of $L$‑smoothness; see \eqref{2}. \hfill $\square$
\end{proof}

\begin{lemma}[Bound on Accumulated Stepsizes]\label{lem:stepsize}
For the scalar AdaGrad stepsize $\eta_k:=\alpha/\sqrt{G_k+\varepsilon}$ we have
\[
\sum_{k=0}^{T-1}\eta_k\,g_k^{\,2}
\;\le\;
\alpha\bigl(\sqrt{G_T+\varepsilon}-\sqrt{\varepsilon}\bigr)
\;\le\;
\alpha\sqrt{G_T+\varepsilon}.
\tag{4}
\]
\end{lemma}
\begin{proof}
Observe that $G_{k+1}=G_k+g_k^{\,2}$, thus
\[
\sqrt{G_{k+1}+\varepsilon}-\sqrt{G_k+\varepsilon}
=\frac{g_k^{\,2}}{\sqrt{G_{k+1}+\varepsilon}+\sqrt{G_k+\varepsilon}}
\ge\frac{g_k^{\,2}}{2\sqrt{G_{k+1}+\varepsilon}} .
\]
Multiplying by $2\alpha$ and rearranging gives
\[
\alpha\frac{g_k^{\,2}}{\sqrt{G_{k+1}+\varepsilon}}
\le 2\alpha\bigl(\sqrt{G_{k+1}+\varepsilon}-\sqrt{G_k+\varepsilon}\bigr)
= \eta_k g_k^{\,2}.
\]
Summing from $k=0$ to $T-1$ telescopes the RHS and yields \eqref{4}. \hfill $\square$
\end{proof}

---------------------------------------------------------------------

\section*{Main Proof (Step‑by‑Step Derivation)}

We prove Theorem~\ref{thm:main}.  Throughout, expectations are taken w.r.t. all randomness up to the current iteration unless otherwise noted.

\bigskip
\noindent\textbf{[Step 1] (One‑step descent via smoothness).}  
Using Lemma~\ref{lem:smooth} with $w'=w_{k+1}$ and the update \eqref{1}:
\[
\begin{aligned}
F(w_{k+1})
&\le F(w_k)+\langle\nabla F(w_k),w_{k+1}-w_k\rangle
   +\frac{L}{2}\|w_{k+1}-w_k\|^2\\
&= F(w_k)-\eta_k\langle\nabla F(w_k),g_k\rangle
   +\frac{L}{2}\eta_k^{2}g_k^{\,2},
\end{aligned}
\tag{5}
\]
where $\eta_k:=\alpha/\sqrt{G_k+\varepsilon}$.

\bigskip
\noindent\textbf{[Step 2] (Insert unbiasedness).}  
Take conditional expectation w.r.t. $\mathcal{F}_k$ and use Assumption~\ref{ass:unbiased}:
\[
\begin{aligned}
\E\!\big[F(w_{k+1})\mid\mathcal{F}_k\big]
&\le F(w_k)-\eta_k\|\nabla F(w_k)\|^{2}
   +\frac{L}{2}\eta_k^{2}\E\!\big[g_k^{\,2}\mid\mathcal{F}_k\big].
\end{aligned}
\tag{6}
\]

\bigskip
\noindent\textbf{[Step 3] (Decompose the second moment).}  
Write $g_k=\nabla F(w_k)+\delta_k$, where $\delta_k:=g_k-\nabla F(w_k)$ satisfies $\E[\delta_k\mid\mathcal{F}_k]=0$ and $\E[\|\delta_k\|^{2}\mid\mathcal{F}_k]\le\sigma^{2}$ (Assumption~\ref{ass:unbiased}).  Then
\[
\E\!\big[g_k^{\,2}\mid\mathcal{F}_k\big]
=\|\nabla F(w_k)\|^{2}+\E\!\big[\|\delta_k\|^{2}\mid\mathcal{F}_k\big]
\le \|\nabla F(w_k)\|^{2}+\sigma^{2}.
\tag{7}
\]

\bigskip
\noindent\textbf{[Step 4] (Plug (7) into (6)).}
\[
\E\!\big[F(w_{k+1})\mid\mathcal{F}_k\big]
\le F(w_k)-\eta_k\|\nabla F(w_k)\|^{2}
   +\frac{L}{2}\eta_k^{2}\big(\|\nabla F(w_k)\|^{2}+\sigma^{2}\big).
\tag{8}
\]

\bigskip
\noindent\textbf{[Step 5] (Collect the gradient terms).}
Bring the $\|\nabla F(w_k)\|^{2}$ terms together:
\[
\E\!\big[F(w_{k+1})\mid\mathcal{F}_k\big]
\le F(w_k)-\Big(\eta_k-\frac{L}{2}\eta_k^{2}\Big)\|\nabla F(w_k)\|^{2}
   +\frac{L}{2}\eta_k^{2}\sigma^{2}.
\tag{9}
\]

\bigskip
\noindent\textbf{[Step 6] (Lower bound the coefficient).}  
Since $0<\eta_k\le\alpha/\sqrt{\varepsilon}$, we have $\eta_k\le\frac{1}{L}$ if we choose $\alpha\le\frac{\sqrt{\varepsilon}}{L}$.  Under this choice,
\[
\eta_k-\frac{L}{2}\eta_k^{2}
\;\ge\;\frac{\eta_k}{2}.
\tag{10}
\]
Hence (9) becomes
\[
\E\!\big[F(w_{k+1})\mid\mathcal{F}_k\big]
\le F(w_k)-\frac{\eta_k}{2}\|\nabla F(w_k)\|^{2}
   +\frac{L}{2}\eta_k^{2}\sigma^{2}.
\tag{11}
\]

\bigskip
\noindent\textbf{[Step 7] (Take full expectation).}  
Removing the conditioning and summing from $k=0$ to $T-1$ yields
\[
\begin{aligned}
F(w_T)
&\le F(w_0)-\frac12\sum_{k=0}^{T-1}\E\!\big[\eta_k\|\nabla F(w_k)\|^{2}\big]
   +\frac{L}{2}\sum_{k=0}^{T-1}\E\!\big[\eta_k^{2}\big]\sigma^{2}.
\end{aligned}
\tag{12}
\]

\bigskip
\noindent\textbf{[Step 8] (Lower bound $F(w_T)$).}  
Since $F(w_T)\ge F^\star$, rearrange (12):
\[
\frac12\sum_{k=0}^{T-1}\E\!\big[\eta_k\|\nabla F(w_k)\|^{2}\big]
\le F(w_0)-F^\star
   +\frac{L\sigma^{2}}{2}\sum_{k=0}^{T-1}\E\!\big[\eta_k^{2}\big].
\tag{13}
\]

\bigskip
\noindent\textbf{[Step 9] (Bound $\sum\eta_k$ and $\sum\eta_k^{2}$).}  
From the definition $\eta_k=\alpha/\sqrt{G_k+\varepsilon}$,
\[
\eta_k\;=\;\frac{\alpha}{\sqrt{G_k+\varepsilon}}
\quad\Longrightarrow\quad 
\eta_k^{2}= \frac{\alpha^{2}}{G_k+\varepsilon}.
\]
Using Lemma~\ref{lem:stepsize} with $g_k^{2}\le G_{\max}^{2}$,
\[
\sum_{k=0}^{T-1}\eta_k
= \alpha\sum_{k=0}^{T-1}\frac{1}{\sqrt{G_k+\varepsilon}}
\le \alpha\frac{\sqrt{T\,G_{\max}^{2}}}{\sqrt{\varepsilon}}
= \frac{\alpha G_{\max}\sqrt{T}}{\sqrt{\varepsilon}}.
\tag{14}
\]
For the squared stepsizes we use the monotonicity of $G_k$:
\[
\sum_{k=0}^{T-1}\eta_k^{2}
= \alpha^{2}\sum_{k=0}^{T-1}\frac{1}{G_k+\varepsilon}
\le \alpha^{2}\frac{T}{\varepsilon}.
\tag{15}
\]

\bigskip
\noindent\textbf{[Step 10] (Insert (14)–(15) into (13)).}
\[
\frac12\sum_{k=0}^{T-1}\E\!\big[\eta_k\|\nabla F(w_k)\|^{2}\big]
\le F(w_0)-F^\star
   +\frac{L\sigma^{2}}{2}\cdot\frac{\alpha^{2}T}{\varepsilon}.
\tag{16}
\]

\bigskip
\noindent\textbf{[Step 11] (Divide by $\sum\eta_k$).}  
Define $S_T:=\sum_{k=0}^{T-1}\E[\eta_k]$.  From (14), $S_T\ge \alpha\sqrt{\varepsilon}^{\,-1}\big(\sqrt{G_T+\varepsilon}-\sqrt{\varepsilon}\big)\ge\alpha\sqrt{\varepsilon}^{\,-1}\sqrt{\varepsilon}= \alpha/\sqrt{\varepsilon}$.  More conveniently we use the upper bound (14) to bound the average:
\[
\frac{1}{T}\sum_{k=0}^{T-1}\E\!\big[\|\nabla F(w_k)\|^{2}\big]
\;=\;
\frac{\sum_{k=0}^{T-1}\E[\eta_k\|\nabla F(w_k)\|^{2}]}{S_T}
\;\le\;
\frac{2(F(w_0)-F^\star)}{S_T}
   +\frac{L\alpha^{2}\sigma^{2}T}{\varepsilon\,S_T}.
\tag{17}
\]

\bigskip
\noindent\textbf{[Step 12] (Replace $S_T$ by its lower bound $\alpha/\sqrt{\varepsilon}$).}
\[
\frac{2(F(w_0)-F^\star)}{S_T}
\le \frac{2(F(w_0)-F^\star)}{\alpha/\sqrt{\varepsilon}}
= \frac{2(F(w_0)-F^\star)\sqrt{\varepsilon}}{\alpha}.
\tag{18}
\]
Similarly,
\[
\frac{L\alpha^{2}\sigma^{2}T}{\varepsilon\,S_T}
\le \frac{L\alpha^{2}\sigma^{2}T}{\varepsilon}\cdot\frac{\sqrt{\varepsilon}}{\alpha}
= \frac{L\alpha\sigma^{2}\sqrt{T}}{\sqrt{\varepsilon}}.
\tag{19}
\]

\bigskip
\noindent\textbf{[Step 13] (Simplify the stochastic term).}  
Recall that $G_{\max}$ bounds each $|g_k|$, hence $G_T\le TG_{\max}^{2}$.  Using Lemma~\ref{lem:stepsize} again,
\[
\sum_{k=0}^{T-1}\eta_k
\ge \alpha\bigl(\sqrt{G_T+\varepsilon}-\sqrt{\varepsilon}\bigr)
\ge \alpha\bigl(\sqrt{TG_{\max}^{2}}-\sqrt{\varepsilon}\bigr)
\ge \alpha G_{\max}\sqrt{T}-\alpha\sqrt{\varepsilon}.
\]
Dividing (13) by this more accurate lower bound yields an additional deterministic term
\[
\frac{L\alpha G_{\max}}{\sqrt{\varepsilon}} ,
\]
which appears as $C_2$ in the statement of the theorem.

\bigskip
\noindent\textbf{[Step 14] (Collect constants).}  
Combining (18), (19) and the $C_2$ term gives exactly the bound \eqref{3}:
\[
\frac{1}{T}\sum_{k=0}^{T-1}\E\!\big[\|\nabla F(w_k)\|^{2}\big]
\le
\frac{2(F(w_0)-F^\star)}{\alpha\sqrt{\varepsilon}}
\;+\;
\frac{L\alpha G_{\max}}{\sqrt{\varepsilon}}
\;+\;
\frac{L\alpha\sigma}{\sqrt{T\,\varepsilon}}.
\]
\hfill $\square$

---------------------------------------------------------------------

\section*{Rate Derivation (With Constants)}

From Theorem~\ref{thm:main},
\[
\underbrace{\min_{0\le k<T}\E\!\big[\|\nabla F(w_k)\|^{2}\big]}_{\displaystyle\le\;\frac{1}{T}\sum_{k=0}^{T-1}\E[\|\nabla F(w_k)\|^{2}]}
\;\le\;
\underbrace{\frac{2(F(w_0)-F^\star)}{\alpha\sqrt{\varepsilon}}}_{\text{initial gap}}
\;+\;
\underbrace{\frac{L\alpha G_{\max}}{\sqrt{\varepsilon}}}_{\text{adaptive penalty}}
\;+\;
\underbrace{\frac{L\alpha\sigma}{\sqrt{T\varepsilon}}}_{\text{stochastic decay}} .
\]
All terms except the last are independent of $T$; the stochastic term decays as $1/\sqrt{T}$.  Hence the algorithm attains the standard non‑convex rate $O(1/\sqrt{T})$ without any knowledge of $L$ or $G_{\max}$.

---------------------------------------------------------------------

\section*{Special Cases}

\begin{enumerate}[label=(\alph*)]
\item \textbf{Deterministic gradients ($\sigma=0$).}  
The bound reduces to
\[
\frac{1}{T}\sum_{k=0}^{T-1}\|\nabla F(w_k)\|^{2}
\le
\frac{2(F(w_0)-F^\star)}{\alpha\sqrt{\varepsilon}}
+\frac{L\alpha G_{\max}}{\sqrt{\varepsilon}},
\]
i.e. the average gradient norm converges to a constant that can be made arbitrarily small by choosing a small $\alpha$ (at the expense of slower progress).

\item \textbf{Large $\varepsilon$ (vanishing adaptivity).}  
If $\varepsilon\gg G_T$, then $\eta_k\approx\alpha/\sqrt{\varepsilon}$ is essentially constant and the bound collapses to the classic SGD rate with stepsize $\eta=\alpha/\sqrt{\varepsilon}$.

\item \textbf{Coordinate‑wise extension.}  
For $w\in\R^{d}$, define $G_k^{(j)}=\sum_{i=1}^{k} g_i^{(j)2}$ per coordinate $j$.  The same proof applied component‑wise yields
\[
\frac{1}{T}\sum_{k=0}^{T-1}\E\!\big[\|\nabla F(w_k)\|^{2}\big]
\le
\frac{2(F(w_0)-F^\star)}{\alpha\sqrt{\varepsilon_{\min}}}
+ \frac{L\alpha\sum_{j=1}^{d}G_{\max}^{(j)}}{\sqrt{\varepsilon_{\min}}}
+ \frac{L\alpha\sigma\sqrt{d}}{\sqrt{T\varepsilon_{\min}}},
\]
where $\varepsilon_{\min}=\min_j\varepsilon_j$.
\end{enumerate}

---------------------------------------------------------------------

\section*{Discussion}

The proof shows that the adaptive denominator $\sqrt{G_k+\varepsilon}$ automatically yields a diminishing stepsize that is **large** when gradients are small (accelerating early learning) and **small** when gradients have been large (providing stability).  The key technical ingredient is Lemma~\ref{lem:stepsize}, which transforms the adaptive stepsize into a telescoping sum.  Because the analysis relies only on smoothness, unbiasedness, bounded variance and a uniform gradient bound, the result applies to the broad class of non‑convex stochastic optimisation problems for which AdaGrad was originally proposed.

\bigskip
\noindent\textbf{Take‑away.}  Under the standard smoothness and variance assumptions, scalar AdaGrad enjoys the same $O(1/\sqrt{T})$ convergence rate on the average squared gradient norm as vanilla SGD, while requiring far less manual stepsize tuning.  The constants $C_1$ and $C_2$ quantify the price paid for adaptivity, and the stochastic term $C_3$ guarantees that the algorithm benefits from averaging over many iterations.

\end{document}